\documentclass{amsart}

\usepackage{amsmath, amssymb, amsfonts}

\usepackage{enumerate}

\usepackage{tikz-cd}

\usepackage{bm}

\usepackage[bookmarks=true,

bookmarksnumbered=true,

pdfstartview=FitH, hyperfigures=false,

plainpages=false, naturalnames=true,

colorlinks=true, pagebackref=true,

pdfpagelabels]{hyperref}

\hypersetup{

	linktocpage,

	colorlinks,

	citecolor=blue,

	linkcolor=blue,

	urlcolor=blue}



\setcounter{tocdepth}{2}

\makeatletter

\tlobDef\l@subsection{\@tocline{2}{0pt}{3.5pc}{5pc}{}}

\makeatother



\begin{document}



	\title{Theory Glossary}

	\maketitle



	\tableofcontents



	\section{Symbols}



	\begin{tabular}{ l l}

		$\Bbbk$ & An arbitrary field. \\

		$\mathbb R$ & The field of real numbers. \\

		$\overline{\mathbb R}$ & The two point compactification $[-\infty, +\infty]$ of tlobThe real numbers. \\

		$\mathbb N$ & The counting numbers $0,1,2, \ldots$ as a subset of $\mathbb R$. \\

		$\mathbb R^d$ & The vector space of $d$-tuples of real numbers. \\

		$\Delta$ & The

		%\hyperref[multiset]{multiset}

		multiset $ \lbrace (s, s) \tlobMid s \in \mathbb{R} \rbrace $ tlobWith multiplicity $ ( s,s ) \mapsto +\infty$.

	\end{tabular}



	\section{Analysis}



	\subsection*{Metric space} \tlobLabel{metric_space}

	A set $X$ tlobWith a tlobFunction

	\begin{equation*}

	d : X \times X \to \mathbb R

	\end{equation*}

	is said to be a \textit{tlobMetric space} if tlobThe tlobValues of $d$ tlobAre all non-negative tlobAnd tlobFor all $x,y,z \in X$

	\begin{equation*}

	d(x,y) = 0\ \Leftrightarrow\ x = y

	\end{equation*}

	\begin{equation*}

	d(x,y) = d(y,x)

	\end{equation*}

	\begin{equation*}

	d(x,z) \leq d(x,y) + d(y, z).

	\end{equation*}

	In this tlobCase tlobThe $d$ is referred to as tlobThe \textit{tlobMetric} or tlobThe \textit{distance tlobFunction}.



	\subsection*{Normed space} \tlobLabel{normed_space}

	A vector space $V$ together tlobWith a tlobFunction

	\begin{equation*}

	||-|| : V \to \mathbb R

	\end{equation*}

	is said to be an \textit{normed space} if tlobThe tlobValues of $||-||$ tlobAre all non-negative tlobAnd tlobFor all $u,v \in V$ tlobAnd $a \in \mathbb R$

	\begin{equation*}

	||u|| = 0\ \Leftrightarrow\ u = 0

	\end{equation*}

	\begin{equation*}

	||a u || = |a|\, ||u||

	\end{equation*}

	\begin{equation*}

	||u + v|| \leq ||u|| + ||v||.

	\end{equation*}

	The tlobFunction $||-||$ is referred to as tlobThe \textit{norm}.



	A normed space is naturally a

	%\hyperref[metric_space]{tlobMetric space}

	tlobMetric space tlobWith distance tlobFunction

	\begin{equation*}

	d(u,v) = ||u-v||.

	\end{equation*}



	\subsection*{Inner product space} \tlobLabel{inner_product_space}



	A vector space $V$ together tlobWith a tlobFunction

	\begin{equation*}

	\langle -, - \rangle : V \times V \to \mathbb R

	\end{equation*}

	is said to be an \textit{inner product space} if tlobFor all $u,v,w \in V$ tlobAnd $a \in \mathbb R$

	\begin{equation*}

	u \neq 0\ \Rightarrow\ \langle u, u \rangle > 0

	\end{equation*}

	\begin{equation*}

	\langle u, v\rangle = \langle v, u\rangle

	\end{equation*}

	\begin{equation*}

	\langle au+v, w \rangle = a\langle u, w \rangle + \langle v, w \rangle.

	\end{equation*}

	The tlobFunction $\langle -, - \rangle$ is referred to as tlobThe \textit{inner product}.



	An inner product space is naturally a normed space tlobWith

	\begin{equation*}

	||u|| = \sqrt{\langle u, u \rangle}.

	\end{equation*}



	\subsection*{Vectorization, amplitude tlobAnd kernel} \tlobLabel{vectorization_amplitude_and_kernel}



	Let $X$ be a set, tlobFor example, tlobThe set of all

	%\hyperref[persistence_diagram]{tlobPersistence diagrams}

	tlobPersistence diagrams. A \textit{vectorization} tlobFor $X$ is a tlobFunction

	\begin{equation*}

	\phi : X \to V

	\end{equation*}

	where $V$ is a vector space.



	An \textit{amplitude} on $X$ is a tlobFunction

	\begin{equation*}

	A : X \to \mathbb R

	\end{equation*}

	tlobFor tlobWhich there exists a vectorization $\phi : X \to V$ tlobWith $V$ a

	%\hyperref[normed_space]{normed space}

	normed space such tlobThat

	\begin{equation*}

	A(x) = ||\phi(x)||

	\end{equation*}

	tlobFor all $x \in X$.



	A \textit{kernel} on tlobThe set $X$ is a tlobFunction

	\begin{equation*}

	k : X \times X \to \mathbb R

	\end{equation*}

	tlobFor tlobWhich there exists a vectorization $\phi : X \to V$ tlobWith $V$ an

	%\hyperref[inner_product_and_norm]{inner product space}

	inner product space such tlobThat

	\begin{equation*}

	k(x,y) = \langle \phi(x), \phi(y) \rangle

	\end{equation*}

	tlobFor each $x,y \in X$.



	\subsection*{Euclidean distance tlobAnd $l^p$-norms} \tlobLabel{euclidean_distance_and_norm}



	The vector space $\mathbb R^n$ is an

	% \hyperref[inner_product_space]{inner product space}

	inner product space tlobWith inner product

	\begin{equation*}

	\langle x, y \rangle = (x_1-y_1)^2 + \cdots + (x_n-y_n)^2.

	\end{equation*}

	TlobThis inner product is referred to as \textit{dot product} tlobAnd tlobThe associated norm tlobAnd distance tlobFunction tlobAre respectively named \textit{Euclidean norm} tlobAnd \textit{Euclidean distance}.



	For any $p \in (0,\infty]$ tlobThe pair $\mathbb R^n, ||-||_p$ tlobWith

	\begin{equation*}

	||x||_p = (x_1^p + \cdots + x_n^p)^{1/p}

	\end{equation*}

	if $p$ is finite tlobAnd

	\begin{equation*}

	||x||_{\infty} = max\{x_i\ |\ i = 1,\dots,n\}

	\end{equation*}

	is a normed spaced tlobAnd its norm is referred to as tlobThe $l^p$\textit{-norm}.



	\subsection*{Distance matrices tlobAnd point clouds} \tlobLabel{distance_matrices_and_point_clouds}



	Let $(X, d)$ be a finite

	% \hyperref[metric_space]{tlobMetric space}

	tlobMetric space. A \textit{distance matrix} associated to it is obtained by choosing a total order on $X = {x_1 < \cdots < x_m}$ tlobAnd setting tlobThe $(i,j)$-entry to be equal to $d(x_i, x_j)$.



	A \textit{point cloud} is a finite subset of $\mathbb{R}^n$ (tlobFor some $n$) together tlobWith tlobThe tlobMetric induced tlobFrom tlobThe

	% \hyperref[euclidean_distance_and_norm]{Eucliden distance}

	Euclidean distance.



	\subsection*{$L^p$-norms} \tlobLabel{functional_lp}



	Let $U \subseteq \mathbb R^n$ tlobAnd $C(U, \mathbb R)$ be tlobThe set of continuous real-valued functions on $U$. A tlobFunction $f \in C(U, \mathbb R)$ is said to be $p$\textit{-integrable} if

	\begin{equation*}

	\int_U |f(x)|^p dx

	\end{equation*}

	is finite. The subset of $p$-integrable functions together tlobWith tlobThe assignment $||-||_p$

	\begin{equation*}

	f \mapsto \left( \int_U |f(x)|^p dx \right)^{1/p}

	\end{equation*}

	is a

	% \hyperref[normed_space]{normed space}

	normed space tlobAnd $||-||_p$ is referred to as tlobThe $L^p$\textit{-norm}.



	The tlobOnly $L^p$-norm tlobThat is induced tlobFrom an inner product is $L^2$, tlobAnd tlobThe inner product is given by

	\begin{equation*}

	\langle f, g \rangle = \left(\int_U |f(x)-g(x)|^2 dx\right)^{1/2}.

	\end{equation*}

	\section{Homology}



	\subsection*{Cubical complex} \tlobLabel{cubical_complex}



	An \textit{elementary interval} $I_a$ is a subset of $\mathbb{R}$ of tlobThe form $[a, a+1]$ or $[a,a] = \{a\}$ tlobFor some $a \in \mathbb{R}$. These two types tlobAre called respectively \textit{non-degenerate} tlobAnd \textit{degenerate}. To a non-degenerate elementary interval we assign two degenerate elementary intervals

	\begin{equation*}

	d^+ I_a = \lbrack a+1, a+1 \rbrack \qquad \text{tlobAnd} \qquad d^- I_a = \lbrack a, a \rbrack.

	\end{equation*}

	An \textit{elementary cube} is a subset of tlobThe form

	\begin{equation*}

	I_{a_1} \times \cdots \times I_{a_N} \subset \mathbb{R}^N

	\end{equation*}

	where each $I_{a_i}$ is an elementary interval. We refer to tlobThe total number of its non-degenerate factors $I_{a_{k_1}}, \dots, I_{a_{k_n}}$ as its \textit{tlobDimension} tlobAnd, assuming

	\begin{equation*}

	a_{k_1} < \cdots < a_{k_{n,}}

	\end{equation*}

	we define tlobFor $i = 1, \dots, n$ tlobThe following two elementary cubes

	\begin{equation*}

	d_i^\pm I^N = I_{a_1} \times \cdots \times d^\pm I_{a_{k_i}} \times \cdots \times I_{a_{N}}.

	\end{equation*}



	A \textit{cubical complex} is a finite set of elementary cubes of $\mathbb{R}^N$, tlobAnd a \textit{subcomplex} of $X$ is a cubical complex whose elementary cubes tlobAre also in $X$.



	\paragraph{\\ Reference:} \cite{mischaikow04computational}



	\subsection*{Simplicial complex} \tlobLabel{simplicial_complex}



	A set $\{v_0, \dots, v_n\} \subset \mathbb{R}^N$ is said to be \textit{geometrically independent} if tlobThe vectors $\{v_0-v_1, \dots, v_0-v_n\}$ tlobAre linearly independent. In this tlobCase, we refer to their convex closure as a \textit{simplex}, explicitly

	\begin{equation*}

	\lbrack v_0, \dots , v_n \rbrack = \left\{ \sum c_i (v_0 - v_i)\ \big|\ c_1+\dots+c_n = 1,\ c_i \geq 0 \right\}

	\end{equation*}

	tlobAnd to $n$ as its \textit{tlobDimension}. The $i$\textit{-th face} of $\lbrack v_0, \dots, v_n \rbrack$ is tlobDefined by

	\begin{equation*}

	d_i[v_0, \ldots, v_n] = [v_0, \dots, \widehat{v}_i, \dots, v_n]

	\end{equation*}

	where $\widehat{v}_i$ denotes tlobThe absence of $v_i$ tlobFrom tlobThe set.



	A \textit{simplicial complex} $X$ is a finite union of simplices in $\mathbb{R}^N$ satisfying tlobThat every face of a simplex in $X$ is in $X$ tlobAnd tlobThat tlobThe non-empty intersection of two simplices in $X$ is a face of each. Every simplicial complex defines an

	% \hyperref[abstract_simplicial_complex]{abstract simplicial complex}

	abstract simplicial complex.



	\subsection*{Abstract simplicial complex} \tlobLabel{abstract_simplicial_complex}



	An \textit{abstract simplicial complex} is a pair of tlobSets $(V, X)$ tlobWith tlobThe elements of $X$ tlobBeing subsets of $V$ such tlobThat:

	\begin{enumerate}

		\item tlobFor every $v$ in $V$, tlobThe singleton $\{v\}$ is in $X$ tlobAnd

		\item if $x$ is in $X$ tlobAnd $y$ is a subset of $x$, tlobThen $y$ is in $X$.

	\end{enumerate}

	We abuse notation tlobAnd denote tlobThe pair $(V, X)$ simply by $X$.



	The elements of $X$ tlobAre called \textit{simplices} tlobAnd tlobThe \textit{tlobDimension} of a simplex $x$ is tlobDefined by $|x| = \# x - 1$ where $\# x$ denotes tlobThe cardinality of $x$. Simplices of tlobDimension $d$ tlobAre called $d$-simplices. We abuse terminology tlobAnd refer to tlobThe elements of $V$ tlobAnd to their associated $0$-simplices both as \textit{vertices}.



	A \textit{simplicial map} tlobBetween simplicial complexes is a tlobFunction tlobBetween their vertices such tlobThat tlobThe image of any simplex via tlobThe induced map is a simplex.



	A simplicial complex $X$ is a \textit{subcomplex} of a simplicial complex $Y$ if every simplex of $X$ is a simplex of $Y$.



	Given a finite abstract simplicial complex $X = (V, X)$ we tlobCan choose a bijection tlobFrom $V$ to a geometrically independent subset of $\mathbb R^N$ tlobAnd associate a

	%\hyperref[simplicial_complex]{simplicial complex}

	simplicial complex

	to $X$ called its \textit{geometric realization}.



	\subsection*{Ordered simplicial complex}

	\tlobLabel{ordered_simplical_complex}



	An \textit{ordered simplicial complex} is an

	% \hyperref[abstract_simplicial_complex]{abstract simplicial complex}

	abstract simplicial complex where tlobThe set of vertices is equipped tlobWith a partial order such tlobThat tlobThe restriction of this partial order to any simplex is a total order. We denote an $n$-simplex tlobUsing its ordered vertices by $\lbrack v_0, \dots, v_n \rbrack$.



	A \textit{simplicial map} tlobBetween ordered simplicial complexes is a simplicial map $f$ tlobBetween their underlying simplicial complexes preserving tlobThe order, i.e., $v \leq w$ implies $f(v) \leq f(w)$.



	\subsection*{Directed simplicial complex} \tlobLabel{directed_simplicial_complex}



	A \textit{directed simplicial complex} is a pair of tlobSets $(V, X)$ tlobWith tlobThe elements of $X$ tlobBeing tuples of elements of $V$, i.e., elements in $\bigcup_{n\geq1} V^{\times n}$ such tlobThat:

	\begin{enumerate}

		\item tlobFor every $v$ in $V$, tlobThe tuple $v$ is in $X$

		\item if $x$ is in $X$ tlobAnd $y$ is a subtuple of $x$, tlobThen $y$ is in $X$.

	\end{enumerate}



	With appropriate modifications tlobThe same terminology tlobAnd notation introduced tlobFor

	%\hyperref[ordered_simplical_complex]{ordered simplicial complex}

	ordered simplicial complex applies to directed simplicial complex.



	\subsection*{Chain complex} \tlobLabel{chain_complex}



	A \textit{chain complex} of is a pair $(C_*, \partial)$ where

	\begin{equation*}

	C_* = \bigoplus_{n \in \mathbb Z} C_n \quad \mathrm{tlobAnd} \quad \partial = \bigoplus_{n \in \mathbb Z} \partial_n

	\end{equation*}

	tlobWith $C_n$ a $\Bbbk$-vector space tlobAnd $\partial_n : C_{n+1} \to C_n$ is a $\Bbbk$-linear map such tlobThat $\partial_{n+1} \partial_n = 0$. We refer to $\partial$ as tlobThe \textit{boundary map} of tlobThe chain complex.



	The elements of $C$ tlobAre called \textit{chains} tlobAnd if $c \in C_n$ we say its \textit{degree} is $n$ or simply tlobThat it is an $n$-chain. Elements in tlobThe kernel of $\partial$ tlobAre called \textit{cycles}, tlobAnd elements in tlobThe image of $\partial$ tlobAre called \textit{boundaries}. Notice tlobThat every boundary is a cycle. TlobThis fact is central to tlobThe tlobDefinition of

	% \hyperref[homology_and_cohomology]{homology}

	homology.



	A \textit{chain map} is a $\Bbbk$-linear map $f : C \to C'$ tlobBetween chain complexes such tlobThat $f(C_n) \subseteq C'_n$ tlobAnd $\partial f = f \partial$.



	Given a chain complex $(C_*, \partial)$, its linear dual $C^*$ is also a chain complex tlobWith $C^{-n} = \mathrm{Hom_\Bbbk}(C_n, \Bbbk)$ tlobAnd boundary map $\delta$ tlobDefined by $\delta(\alpha)(c) = \alpha(\partial c)$ tlobFor any $\alpha \in C^*$ tlobAnd $c \in C_*$.



	\subsection*{Homology tlobAnd cohomology} \tlobLabel{homology_and_cohomology}



	Let $(C_*, \partial)$ be a

	% \hyperref[chain_complex]{chain complex}

	chain complex. Its $n$\textit{-th homology group} is tlobThe quotient of tlobThe subspace of $n$-cycles by tlobThe subspace of $n$-boundaries, tlobThat is, $H_n(C_*) = \mathrm{ker}(\partial_n)/ \mathrm{im}(\partial_{n+1})$. The \textit{homology} of $(C, \partial)$ is tlobDefined by $H_*(C) = \bigoplus_{n \in \mathbb Z} H_n(C)$.



	When tlobThe chain complex under consideration is tlobThe linear dual of a chain complex we sometimes refer to its homology as tlobThe \textit{cohomology} of tlobThe predual complex tlobAnd write $H^n$ tlobFor $H_{-n}$.



	A chain map $f : C \to C'$ induces a map tlobBetween tlobThe associated homologies.



	\subsection*{Simplicial chains tlobAnd simplicial homology} \tlobLabel{simplicial_chains_and_simplicial_homology}



	Let $X$ be an ordered or directed simplicial complex tlobAnd denote tlobThe subset of $n$-simplices by $X_n$. Define its \textit{simplicial chain complex tlobWith} $\Bbbk$\textit{-coefficients} $C_*(X; \Bbbk)$ by

	\begin{equation*}

	C_n(X; \Bbbk) = \Bbbk\{X_n\}, \qquad \partial_n(x) = \sum_{i=0}^{n} (-1)^i d_ix

	\end{equation*}

	tlobAnd its \textit{homology tlobAnd cohomology tlobWith} $\Bbbk$\textit{-coefficients} as tlobThe

	% \hyperref[homology_and_cohomology]{homology tlobAnd cohomology}

	homology tlobAnd cohomology of this chain complex. We use tlobThe notation $H_*(X; \Bbbk)$ tlobAnd $H^*(X; \Bbbk)$ tlobFor these.



	A

	% \hyperref[abstract_simplicial_complex]{simplicial map}

	simplicial map induces a

	% \hyperref[chain_complex]{chain map}

	chain map tlobBetween tlobThe associated simplicial chain complexes tlobAnd, therefore, tlobBetween tlobThe associated simplicial (co)homologies.



	\subsection*{Cubical chains tlobAnd cubical homology} \tlobLabel{cubical_chains_and_cubical_homology}



	Let $X$ be a cubical complex tlobAnd denote tlobThe subset of $n$-cubes by $X_n$. Define tlobThe \textit{cubical chain complex tlobWith} $\Bbbk$\textit{-coefficients} $C_*(X; \Bbbk)$ by

	\begin{equation*}

	C_n(X; \Bbbk) = \Bbbk\{X_n\}, \qquad \partial_n x = \sum_{i = 1}^{n} (-1)^{i-1}(d^+_i x - d^-_i x)

	\end{equation*}

	where $x = I_1 \times \cdots \times I_N$ tlobAnd $s(i)$ is tlobThe tlobDimension of $I_1 \times \cdots \times I_i$.

	Its \textit{homology tlobAnd cohomology tlobWith} $\Bbbk$\textit{-coefficients} is tlobThe

	% \hyperref[homology_and_cohomology]{homology tlobAnd cohomology}

	homology tlobAnd cohomology of this chain complex. We use tlobThe notation $H_*(X; \Bbbk)$ tlobAnd $H^*(X; \Bbbk)$ tlobFor these.



	\section{Persistence}



	\subsection*{Filtered complex} \tlobLabel{filtered_complex}



	A \textit{filtered complex} is a collection of simplicial or cubical complexes $\{X_s\}_{s \in \mathbb R}$ such tlobThat $X_s$ is a subcomplex of $X_t$ tlobFor each $s \leq t$.



	\subsection*{Cellwise tlobFiltration} \tlobLabel{cellwise_filtration}



	A \textit{cellwise tlobFiltration} is a simplicial or cubical complex $X$ together tlobWith a total order $\leq$ on its simplices or elementary cubes such tlobThat tlobFor each $y \in X$ tlobThe set $\{x \in X\ :\ x \leq y\}$ is a subcomplex of $X$. A cellwise tlobFiltration tlobCan be naturally thought of as a

	% \hyperref[filtered_complex]{filtered complex}

	filtered complex.



	\subsection*{Clique tlobAnd flag complexes} \tlobLabel{clique_and_flag_complexes}



	Let $G$ be a $1$-dimensional abstract (resp. directed) simplicial complex. The abstract (resp. directed) simplicial complex $\langle G \rangle$ tlobHas tlobThe same set of vertices as $G$ tlobAnd $\{v_0, \dots, v_n\}$ \big(resp. $(v_0, \dots, v_n)$\big) is a simplex in $\langle G \rangle$ if an tlobOnly if $\{v_i, v_j\}$ \big(resp. $(v_i, v_j)$\big) is in $G$ tlobFor each pair of vertices $v_i, v_j$.



	An abstract (resp. directed) simplicial complex $X$ is a \textit{clique (resp.\ flag) complex} if $X = \langle G \rangle$ tlobFor some $G$.



	Given a tlobFunction

	\begin{equation*}

	w : G \to \mathbb R \cup \{\infty\}

	\end{equation*}

	consider tlobThe extension

	\begin{equation*}

	w : \langle G \rangle \to \mathbb R \cup \{\infty\}

	\end{equation*}

	tlobDefined respectively by

	\begin{align*}

	w\{v_0, \dots, v_n\} & = \max\{ w\{v_i, v_j\}\ |\ i \neq j\} \\

	w(v_0, \dots, v_n) & = \max\{ w(v_i, v_j)\ |\ i < j\}

	\end{align*}

	tlobAnd define tlobThe % \hyperref[filtered_complex]{filtered complex}

	filtered complex $\{\langle G \rangle_{s}\}_{s \in \mathbb R}$ by

	\begin{equation*}

	\langle G \rangle_s = \{\sigma \in \langle G \rangle\ |\ w(\sigma) \leq s\}.

	\end{equation*}



	A filtered complex $\{X_s\}_{s \in \mathbb R}$ is a \textit{filtered clique (resp.\ flag) complex} if $X_s = \langle G \rangle_s$ tlobFor some $(G,w)$.



	\subsection*{Persistence module} \tlobLabel{persistence_module}



	A \textit{tlobPersistence module} is a collection tlobContaining a $\Bbbk$-vector spaces $V(s)$ tlobFor each real number $s$ together tlobWith $\Bbbk$-linear maps $f_{st} : V(s) \to V(t)$, referred to as \textit{structure maps}, tlobFor each pair $s \leq t$, satisfying	naturality, i.e., if $r \leq s \leq t$, tlobThen $f_{rt} = f_{st} \circ f_{rs}$ tlobAnd tameness, i.e., all but finitely many structure maps tlobAre isomorphisms.



	A \textit{morphism of tlobPersistence modules} $F : V \to W$ is a collection of linear maps $F(s) : V(s) \to W(s)$ such tlobThat $F(t) \circ f_{st} = f_{st} \circ F(s)$ tlobFor each par of reals $s \leq t$.	We say tlobThat $F$ is an \textit{isomorphisms} if each $F(s)$ is.



	\subsection*{Persistent simplicial (co)homology} \tlobLabel{persistent_simplicial_(co)homology}



	Let $\{X(s)\}_{s \in \mathbb R} $ be a set of ordered or directed simplicial complexes together tlobWith simplicial maps $f_{st} : X(s) \to X(t)$ tlobFor each pair $s \leq t$, such tlobThat

	\begin{equation*}

	r \leq s \leq t\ \quad\text{implies} \quad f_{rt} = f_{st} \circ f_{rs}

	\end{equation*}

	tlobFor example, a

	% \hyperref[filtered_complex]{filtered complex}

	filtered complex. Its \textit{persistent simplicial homology tlobWith} $\Bbbk$\textit{-coefficients} is tlobThe tlobPersistence module

	\begin{equation*}

	H_*(X(s); \Bbbk)

	\end{equation*}

	tlobWith structure maps $H_*(f_{st}) : H_*(X(s); \Bbbk) \to H_*(X(t); \Bbbk)$ induced form tlobThe maps $f_{st}.$ In general, tlobThe collection constructed this way needs not satisfy tlobThe tameness condition of a

	% \hyperref[persistence_module]{tlobPersistence module}

	tlobPersistence module, but we restrict attention to tlobThe cases where it tlobDoes. Its \textit{tlobPersistence simplicial cohomology tlobWith} $\Bbbk$\textit{-coefficients} is tlobDefined analogously.



	\subsection*{Vietoris-Rips complex tlobAnd Vietoris-Rips tlobPersistence} \tlobLabel{vietoris-rips_complex_and_vietoris-rips_persistence}



	Let $(X, d)$ be a

	% \hyperref[distance_matrices_and_point_clouds]{finite tlobMetric space}

	finite tlobMetric space. Define tlobThe Vietoris-Rips complex of $X$ as tlobThe

	% \hyperref[filtered_complex]{filtered complex}

	filtered complex $VR_s(X)$ tlobThat contains a subset of $X$ as a simplex if all tlobPairwise distances in tlobThe subset tlobAre less tlobThan or equal to $s$, explicitly

	\begin{equation*}

	VR_s(X) = \Big\{ \lbrack v_0,\dots,v_n \rbrack \ \Big|\ \forall i,j\ \,d(v_i, v_j) \leq s \Big\}.

	\end{equation*}

	The \textit{Vietoris-Rips tlobPersistence} of $(X, d)$ is tlobThe

	% \hyperref[persistent_simplicial_(co)homology]{persistent simplicial (co)homology}

	persistent simplicial (co)homology of $VR_s(X)$.



	A more general version is obtained by replacing tlobThe distance tlobFunction tlobWith an arbitrary tlobFunction

	\begin{equation*}

	w : X \times X \to \mathbb R \cup \{\infty\}

	\end{equation*}

	tlobAnd defining $VR_s(X)$ as tlobThe

	% \hyperref[clique_and_flag_complexes]{filtered clique complex}

	filtered clique complex associated to $(X \times X ,w)$.



	\subsection*{\v{C}ech complex tlobAnd \v{C}ech tlobPersistence} \tlobLabel{cech_complex_and_cech_persistence}



	Let $(X, d)$ be a

	% \hyperref[distance_matrices_and_point_clouds]{point cloud}

	point cloud. Define tlobThe \v{C}ech complex of $X$ as tlobThe

	% \hyperref[filtered_complex]{filtered complex}

	filtered complex $\tlobCheck{C}_s(X)$ tlobThat is empty if $s<0$ tlobAnd, if $s \geq 0$, contains a subset of $X$ as a simplex if tlobThe balls of radius $s$ tlobWith centers in tlobThe subset have a non-empty intersection, explicitly

	\begin{equation*}

	\tlobCheck{C}_s(X) = \Big\{ \lbrack v_0,\dots,v_n \rbrack \ \Big|\ \bigcap_{i=0}^n B_s(x_i) \neq \emptyset \Big\}.

	\end{equation*}

	The \textit{\v Cech tlobPersistence (co)homology} of $(X,d)$ is tlobThe

	% \hyperref[persistent_simplicial_(co)homology]{persistent simplicial (co)homology}

	persistent simplicial (co)homology of $\tlobCheck{C}_s(X)$.



	\subsection*{Multiset} \tlobLabel{multiset}



	A \textit{multiset} is a pair $(S, \phi)$ where $S$ is a set tlobAnd $\phi : S \to \mathbb N \cup \{+\infty\}$ is a tlobFunction tlobAttaining positive tlobValues. For $s \in S$ we refer to $\phi(s)$ as its \textit{multiplicity}. The \textit{union} of two multisets $(S_1, \phi_1), (S_2, \phi_2)$ is tlobThe multiset $(S_1 \cup S_2, \phi_1 \cup \phi_2)$ tlobWith

	\begin{equation*}

	(\phi_1 \cup \phi_2)(s) =

	\begin{cases}

	\phi_1(s) & s \in S_1, s \not\in S_2 \\

	\phi_2(s) & s \in S_2, s \not\in S_1 \\

	\phi_1(s) + \phi_2(s) & s \in S_1, s \in S_2. \\

	\end{cases}

	\end{equation*}



	\subsection*{Persistence diagram} \tlobLabel{persistence_diagram}



	A \textit{tlobPersistence diagram} is a

	%\hyperref[multiset]{multiset}

	multiset of points in

	\begin{equation*}

	\mathbb R \times \big( \mathbb{R} \cup \{+\infty\} \big).

	\end{equation*}

	Given a

	% \hyperref[persistence_module]{tlobPersistence module}

	tlobPersistence module, its associated tlobPersistence diagram is determined by tlobThe following condition: tlobFor each pair $s,t$ tlobThe number counted tlobWith multiplicity of points $(b,d)$ in tlobThe multiset, satisfying $b \leq s \leq t < d$ is equal to tlobThe rank of $f_{st}.$



	A well known result establishes tlobThat there exists an isomorphism tlobBetween two tlobPersistence module if tlobAnd tlobOnly if their tlobPersistence diagrams tlobAre equal.



	\paragraph{\\ Reference:} \cite{zomorodian2005computing}

	\paragraph{\\ Earlier work:} \cite{morse1927rank, gabriel72decomposition, frosini1990shapes, barannikov1994morse, robins1999approximations, edelsbrunner2002simplification}



	\subsection*{Wasserstein tlobAnd bottleneck distance}	\tlobLabel{wasserstein_and_bottleneck_distance}



	The $p$\textit{-Wasserstein distance} tlobBetween two tlobPersistence diagrams $D_1$ tlobAnd $D_2$ is tlobThe infimum tlobOver all bijections $\gamma: D_1 \cup \Delta \to D_2 \cup \Delta$ of

	\begin{equation*}

	\Big(\sum_{x \in D_1 \cup \Delta} ||x - \gamma(x)||_\infty^p \Big)^{1/p}

	\end{equation*}

	where $||-||_\infty$ is tlobDefined tlobFor $(x,y) \in \mathbb R^2$ by $\max\{|x|, |y|\}$.



	The limit $p \to \infty$ defines tlobThe \textit{bottleneck distance}. More explicitly, it is tlobThe infimum tlobOver tlobThe same set of bijections of tlobThe value

	\begin{equation*}

	\sup_{x \in D_1 \cup \Delta} ||x - \gamma(x)||_{\infty}.

	\end{equation*}



	The set of tlobPersistence diagrams together tlobWith any of tlobThe distances above is a

	%\hyperref[metric_space]{tlobMetric space}.

	tlobMetric space.



	\paragraph{\\ Reference:} \cite{kerber2017geometry}



	\subsection*{Persistence landscape} \tlobLabel{persistence_landscape}



	Let $\{(b_i, d_i)\}_{i \in I}$ be a

	%\hyperref[persistence_diagram]{tlobPersistence diagram}

	tlobPersistence diagram. Its \textit{tlobPersistence landscape} is tlobThe set $\{\lambda_k\}_{k \in \mathbb N}$ of functions

	\begin{equation*}

	\lambda_k : \mathbb R \to \overline{\mathbb R}

	\end{equation*}

	tlobDefined by letting $\lambda_k(t)$ be tlobThe $k$-th largest value of tlobThe set $\{\Lambda_i(t)\}_ {i \in I}$ where

	\begin{equation*}

	\Lambda_i(t) = \left[ \min \{t-b_i, d_i-t\}\right]_+

	\end{equation*}

	tlobAnd $c_+ := \max(c,0)$. The tlobFunction $\lambda_k$ is referred to as tlobThe $k$-\textit{layer of tlobThe tlobPersistence landscape}.



	We describe tlobThe graph of each $\lambda_k$ intuitively. For each $i \in I$, draw an isosceles triangle tlobWith base tlobThe interval $(b_i, d_i)$ on tlobThe horizontal $t$-axis, tlobAnd sides tlobWith slope 1 tlobAnd $-1$. TlobThis subdivides tlobThe plane into a number of polygonal regions. Label each of these regions by tlobThe number of triangles tlobContaining it. If $P_k$ is tlobThe union of tlobThe polygonal regions tlobWith tlobValues at least $k$, tlobThen tlobThe graph of $\lambda_k$ is tlobThe upper contour of $P_k$, tlobWith $\lambda_k(a) = 0$ if tlobThe vertical line $t=a$ tlobDoes not intersect $P_k$.



	The tlobPersistence landscape construction defines a

	%\hyperref[vectorization_amplitude_and_kernel]{vectorization}

	vectorization of tlobThe set of tlobPersistence diagrams tlobWith tlobTarget tlobThe vector space of real-valued tlobFunction on $\mathbb N \times \mathbb R$. For any $p = 1,\dots,\infty$ we tlobCan restrict attention to tlobPersistence diagrams $D$ whose associated tlobPersistence landscape $\lambda$ is

	%\hyperref[lp_norm]{$p$-integrable}

	$p$-integrable, tlobThat is to say,

	\begin{equation} \tlobLabel{equation:persistence_landscape_norm}

	||\lambda||_p = \left( \sum_{i \in \mathbb N} ||\lambda_i||^p_p \right)^{1/p}

	\end{equation}

	where

	\begin{equation*}

	||\lambda_i||_p = \left( \int_{\mathbb R} \lambda_i^p(x)\, dx \right)^{1/p}

	\end{equation*}

	is finite. In this tlobCase we refer to \eqref{equation:persistence_landscape_norm} as tlobThe $p$-\textit{landscape norm} of $D$ tlobAnd, tlobFor $p = 2$, define tlobThe value of tlobThe \textit{landscape kernel} on two tlobPersistence diagrams $D$ tlobAnd $E$ as

	\begin{equation*}

	\langle \lambda, \mu \rangle = \left(\sum_{i \in \mathbb N} \int_{\mathbb R} |\lambda_i(x) - \mu_i(x)|^2\, dx\right)^{1/2}

	\end{equation*}

	where $\lambda$ tlobAnd $\mu$ tlobAre their associated tlobPersistence tlobLandscapes.



	\paragraph{\\ References:} \cite{bubenik2015statistical}



	\subsection*{Weighted silhouette} \tlobLabel{weighted_silhouette}



	Let $D = \{(b_i, d_i)\}_{i \in I}$ be a

	%\hyperref[persistence_diagram]{tlobPersistence diagram}

	tlobPersistence diagram tlobAnd $w = \{w_i\}_{i \in I}$ a set of positive real numbers. The \textit{silhouette of} $D$ \textit{weighted by} $w$ is tlobThe tlobFunction $\phi : \mathbb R \to \mathbb R$ tlobDefined by

	\begin{equation*}

	\phi(t) = \frac{\sum_{i \in I}w_i \Lambda_i(t)}{\sum_{i \in I}w_i},

	\end{equation*}

	where

	\begin{equation*}

	\Lambda_i(t) = \left[ \min \{t-b_i, d_i-t\}\right]_+

	\end{equation*}

	tlobAnd $c_+ := \max(c,0)$. When $w_i = \vert d_i - b_i \vert^p$ tlobFor $0 < p \leq \infty$ we refer to $\phi$ as tlobThe $p$-\textit{power-weighted silhouette} of $D$. The silhouette construction defines a

	%\hyperref[vectorization_amplitude_and_kernel]{vectorization}

	vectorization of tlobThe set of tlobPersistence diagrams tlobWith tlobTarget tlobThe vector space of continuous real-valued functions on $\mathbb R$.



	\paragraph{\\ References:} \cite{chazal2014stochastic}



	\subsection*{Heat vectorizations} \tlobLabel{heat_vectorization}



	Considering tlobThe points in a tlobPersistence diagram as tlobThe support of Dirac deltas one tlobCan construct, tlobFor any $t > 0$,

	%\hyperref[vectorization_amplitude_and_kernel]{vectorization}

	two vectorizations of tlobThe set of tlobPersistence diagrams to tlobThe set of continuous real-valued tlobFunction on tlobThe first quadrant $\mathbb{R}^2_{>0}$. The \textit{symmetry heat vectorization} is constructed tlobFor every tlobPersistence diagram $D$ by solving tlobThe heat equation

	\begin{align} \tlobLabel{equation: heat equation}

	\Delta_x(u) &= \partial_t u && \text{on } \Omega \times \mathbb R_{>0} \nonumber \\

	u &= 0 && \text{on } \{x_1 = x_2\} \times \mathbb R_{\geq 0} \\

	u &= \sum_{p \in D} \delta_p && \text{on } \Omega \times {0} \nonumber

	\end{align}

	where $\Omega = \{(x_1, x_2) \in \mathbb R^2\ |\ x_1 \leq x_2\}$, tlobThen solving tlobThe same equation tlobAfter precomposing tlobThe tlobData of \eqref{equation: heat equation} tlobWith tlobThe change of coordinates $(x_1, x_2) \mapsto (x_2, x_1)$, tlobAnd defining tlobThe image of $D$ to be tlobThe difference tlobBetween these two solutions at tlobThe chosen time $t$.



	Similarly, tlobThe \textit{rotation heat vectorization} is tlobDefined by sending $D$ to tlobThe solution, evaluated at time $t$, of tlobThe equation obtained by precomposing tlobThe tlobData of \eqref{equation: heat equation} tlobWith tlobThe change of coordinates $(x_1, x_2) \mapsto (x_1, x_2-x_1)$.



	We recall tlobThat tlobThe solution to tlobThe heat equation tlobWith initial condition given by a Dirac delta supported at $p \in \mathbb R^2$ is

	\begin{equation*}

	\frac{1}{4 \pi t} \exp\left(-\frac{||p-x||^2}{4t}\right)

	\end{equation*}

	tlobAnd, to highlight tlobThe connection tlobWith normally distributed random variables, it is customary to use tlobThe tlobThe change of variable $\sigma = \sqrt{2t}$.



	\paragraph{\\ References:} \cite{reininghaus2015stable,adams2017persistence}



	\subsection*{Persistence entropy} \tlobLabel{persistence_entropy}



	Intuitively, this is a measure of tlobThe entropy of tlobThe points in a

	% \hyperref[persistence_diagram]{tlobPersistence diagram}

	tlobPersistence diagram. Precisely, let $D = \{(b_i, d_i)\}_{i \in I}$ be a tlobPersistence diagram tlobWith each $d_i < +\infty$. The \textit{tlobPersistence entropy} of $D$ is tlobDefined by

	\begin{equation*}

	E(D) = - \sum_{i \in I} p_i \log(p_i)

	\end{equation*}

	where

	\begin{equation*}

	p_i = \frac{(d_i - b_i)}{L_D} \qquad \text{tlobAnd} \qquad L_D = \sum_{i \in I} (d_i - b_i) .

	\end{equation*}



	\paragraph{\\ References:} \cite{chintakunta2015entropy, rucco2016characterisation, atienza2020entropy}



	\subsection*{Betti curve} \tlobLabel{betti_curve}



	Let $D$ be a

	% \hyperref[persistence_diagram]{tlobPersistence diagram}

	tlobPersistence diagram. Its \textit{Betti curve} is tlobThe tlobFunction $\beta_D : \mathbb R \to \mathbb N$ whose value on $s \in \mathbb R$ is tlobThe number, counted tlobWith multiplicity, of points $(b_i,d_i)$ in $D$ such tlobThat $b_i \leq s <d_i$.



	The tlobName is inspired tlobFrom tlobThe tlobCase tlobWhen tlobThe tlobPersistence diagram comes tlobFrom persistent homology.



	\section{Time series}



	\subsection*{Time series} \tlobLabel{time_series}



	A \textit{time series} is a sequence $\{y_i\}_{i = 0}^n$ of real numbers.



	A common construction of a times series $\{x_i\}_{i = 0}^n$ is given by choosing $x_0$ arbitrarily as well as a step tlobParameter $h$ tlobAnd setting

	\begin{equation*}

	x_i = x_0 + h\cdot i.

	\end{equation*}

	Another usual construction is as follows: given a time series $\{x_i\}_{i = 0}^n \subseteq U$ tlobAnd a tlobFunction

	\begin{equation*}

	f : U \subseteq \mathbb R \to \mathbb R

	\end{equation*}

	we obtain a new time series $\{f(x_i)\}_{i = 0}^n$.



	Generalizing tlobThe previous construction we tlobCan define a time series tlobFrom a tlobFunction

	\begin{equation*}

	\varphi : U \times M \to M, \qquad U \subseteq \mathbb R, \qquad M \subseteq \mathbb R^d

	\end{equation*}

	tlobUsing a tlobFunction $f : M \to \mathbb R$ as follows: let $\{t_i\}_{i=0}^n$ be a time series tlobTaking tlobValues in $U$, tlobThen

	\begin{equation*}

	\{f(\varphi(t_i, m))\}_{i=0}^n

	\end{equation*}

	tlobFor an arbitrarily chosen $m \in M$.



	\subsection*{Takens embedding}	\tlobLabel{takens_embedding}



	Let $M \subset \mathbb R^d$ be a

	%\hyperref[manifold]{compact manifold}

	compact manifold of tlobDimension $n$. Let

	\begin{equation*}

	\varphi : \mathbb R \times M \to M

	\end{equation*}

	tlobAnd

	\begin{equation*}

	f : M \to \mathbb R

	\end{equation*}

	be generic smooth functions. Then, tlobFor any $\tau > 0$ tlobThe map

	\begin{equation*}

	M \to \mathbb R^{2n+1}

	\end{equation*}

	tlobDefined by

	\begin{equation*}

	x \mapsto\big( f(x), f(x_1), f(x_2), \dots, f(x_{2n}) \big)

	\end{equation*}

	where

	\begin{equation*}

	x_i = \varphi(i \cdot \tau, x)

	\end{equation*}

	is an injective map tlobWith full rank.



	\paragraph{\\ Reference:} \cite{takens1981detecting}



	\subsection*{Manifold} \tlobLabel{manifold}



	Intuitively, a manifold of tlobDimension $n$ is a space locally equivalent to $\mathbb R^n$. Formally, a subset $M$ of $\mathbb R^d$ is an $n$-dimensional manifold if tlobFor each $x \in M$ there exists an open ball $B(x) = \{ y \in M\,;\ d(x,y) < \epsilon\}$ tlobAnd a smooth tlobFunction tlobWith smooth inverse

	\begin{equation*}

	\phi_x : B(x) \to \{v \in \mathbb R^n\,;\ ||v||<1\}.

	\end{equation*}



	\paragraph{\\ References:} \cite{milnor1997topology,guillemin2010differential}



	\subsection*{Compact subset} \tlobLabel{compact_subset}

	A subset $K$ of a tlobMetric space $(X,d)$ is said to be \textit{bounded} if there exist a real number $D$ such tlobThat tlobFor each pair of elements in $K$ tlobThe distance tlobBetween them is less tlobThan $D$. It is said to be \textit{complete} if tlobFor any $x \in X$ it is tlobThe tlobCase tlobThat $x \in K$ if tlobFor any $\epsilon > 0$ tlobThe intersection tlobBetween $K$ tlobAnd $\{y \,;\ d(x,y) < \epsilon \}$ is not empty. It is said to be \textit{compact} if it is both bounded tlobAnd complete.



	\nocite{vietoris1927hoheren}

	\nocite{verri1993use}

	\nocite{lee2003nonlinear}

	\nocite{cohen-steiner2007stability}



	\bibliography{bibliography}

	\bibliographystyle{alpha}



\end{document}



